\section{Transitions and local finiteness of cyclic centers}
\label{sec:transitions}

Continue to assume that a weak measurable five-coloring is given, with
the notation and null-set conventions of Section~\ref{sec:palettes}.
Write $\mathcal I$ for the five-element label set.  As before, $H_i$ is the
closed spatial essential support of $A_i$, and $P_x$ is defined at every
center $x$, modulo angular null sets.  We first use transitions between two colors to obtain circle traces
that avoid both labels.  The resulting angular restrictions show that
centers whose transition graphs contain a cycle form a locally finite
set.  In Section~\ref{sec:interfaces} we construct a continuous
graph-valued map on a square after removing small disks about these
centers.

\subsection{Transition graphs and avoiding traces}

For distinct $i,j\in\mathcal I$, a bounded ball $B\subset\R^2$, and
$h\in\R^2$, define
\[
 D_{ij}(B,h)=\bigl|\{q\in B:
    (c(q),c(q+h))\in\{(i,j),(j,i)\}\}\bigr|,
\]
where vertical bars on a plane set denote Lebesgue area.  The second
endpoint $q+h$ is not required to belong to $B$.  The normalization by
$|h|$ below is chosen to match the circle-averaging estimate used in
Lemma~\ref{lem:transitions-avoiding-trace}: transition sets of area
comparable to $|h|$ will yield traces avoiding both endpoint labels.

\begin{definition}[Transition graph]
\label{def:transitions-graph}
For $x\in\R^2$, let $\Gamma_x$ be the simple graph on $\mathcal I$ in
which $ij$ is an edge precisely when, for every ball $B$ centered at $x$
with positive finite radius,
\begin{equation}
 \limsup_{\substack{h\to0\\h\ne0}}
       \frac{D_{ij}(B,h)}{|h|}>0.
 \label{eq:source-14}
\end{equation}
\end{definition}

For each fixed pair $ij$, its edge-presence locus is closed.  Indeed, a
ball centered at a limit of points of that locus contains a smaller
ball centered at one of those points, and $D_{ij}$ is monotone in its
ball argument.  Moreover,
\[
 ij\in\Gamma_x\quad\Longrightarrow\quad x\in H_i\cap H_j.
\]
To check this in any neighborhood of $x$, use a smaller ball and then
choose $h$ small enough that both endpoints of every transition under
consideration remain in the neighborhood.  A positive-area transition
set gives positive area of each endpoint color there; translation
preserves area.

\begin{lemma}[An edge has an avoiding trace]
\label{lem:transitions-avoiding-trace}
If $ij\in\Gamma_x$, there is a trace $b^{ij}\in\mathcal W_x$ with
$b^{ij}_i=b^{ij}_j=0$ almost everywhere.  Consequently,
\begin{equation}
 P_x(e)\not\subset\{i,j\}
 \quad\text{for almost every }e\in S,
 \qquad ij\in\Gamma_x.
 \label{eq:source-15}
\end{equation}
\end{lemma}

\begin{proof}
Write
\[
 \mathcal Sg(q)=\int_S g(q+e)\,d\sigma(e).
\]
The estimate we need, for every fixed bounded ball $B$ and every label
$s$, is
\begin{equation}
 \bigl\|\mathcal S\bigl(a_s-a_s(\,\cdot+h)\bigr)
       \bigr\|_{L^2(B)}=o\bigl(\sqrt{|h|}\bigr)
 \qquad(h\to0).
 \label{eq:source-16}
\end{equation}
Here and in \eqref{eq:source-14}, the vector $h$ may approach zero in any
direction.

Choose a boundedly supported $g\in L^2(\R^2)$ agreeing with $a_s$ on a
ball containing $B+S+h$ for all sufficiently small $h$, as well as
$B+S$.  The Fourier transform of normalized circle measure satisfies
\[
 |\widehat\sigma(\xi)|\le C(1+|\xi|)^{-1/2}.
\]
For completeness, after rotating coordinates the relevant integral is
a constant multiple of $\int_0^{2\pi}\exp(ir\cos t)\,dt$, where
$r$ is proportional to $|\xi|$.  For $r\ge1$, removing the arcs on
which $|\sin t|\le r^{-1/2}$ costs $O(r^{-1/2})$.  On each remaining
interval, integration by parts using the derivative $-ir\sin t$ gives
endpoint terms $O(r^{-1/2})$, while the integral of
$|\cos t|/(r\sin^2t)$ is also $O(r^{-1/2})$.  The bound for $r\le1$
is immediate.

By Plancherel, the square of the global $L^2$ norm of
$\mathcal S(g-g(\,\cdot+h))$, divided by $|h|$, is bounded by the
integral of $|\widehat g(\xi)|^2$ times
\[
 \frac{C}{1+|\xi|}
       \frac{\min\{C|h|^2|\xi|^2,4\}}{|h|}.
\]
This multiplier is uniformly bounded: split at $|\xi|=|h|^{-1}$.
For each fixed $\xi$ it tends to zero as $h\to0$.  Dominated convergence
therefore proves \eqref{eq:source-16}; the cutoff is fixed before $h$
tends to zero.

For a fixed displacement $h$, almost every root $q$ has both $q,q+h$
typical, and both exact unit-circle samplings centered at these points
typical on their circles.  At a transition from color $s$ to color $t$,
the first circle avoids $s$ and the second avoids $t$, by properness of
the coloring on typical points.  Thus the mass of color $t$ on the
first circle equals
$\mathcal S(a_t-a_t(\,\cdot+h))(q)$.  For a transition between $i$ and
$j$ in either order, the first-circle mass of these two colors together
is consequently bounded by
\[
 \bigl|\mathcal S(a_i-a_i(\,\cdot+h))(q)\bigr|
 +\bigl|\mathcal S(a_j-a_j(\,\cdot+h))(q)\bigr|.
\]

Fix a ball $B$ about $x$.  Since $ij\in\Gamma_x$, there is a constant
$c_B>0$ and arbitrarily small nonzero $h$ for which the transition set
in $B$ has area at least $c_B|h|$.  Cauchy--Schwarz and
\eqref{eq:source-16} show that the average of the preceding mass over
this transition set is at most
\[
 \frac{o(\sqrt{|h|})}{\sqrt{c_B|h|}}=o(1).
\]
Now take balls $B_m$ of radii tending to zero about $x$.  After fixing
$B_m$ and its constant $c_{B_m}$, choose $h_m$ sufficiently small and
then a transition root $q_m\in B_m$ from the preceding full-measure set,
whose first-circle mass of $i,j$ is less than $1/m$.  This order of choice requires no lower bound
on $c_{B_m}$ uniform in $m$.  The circle samples at $q_m$ have radius
one; a weak-* subsequential limit belongs to $\mathcal W_x$ and has
zero mass in the two specified coordinates.  Since it is nonnegative,
sums to one, and is supported within $P_x$ almost everywhere, it also
proves \eqref{eq:source-15}.
\end{proof}

\subsection{Two angular state lemmas}

A \emph{state} at a center $x$ will mean a measurable angular set
together with a specified nonempty group of labels contained in
$P_x(e)$ almost everywhere on that set.  For a measurable set
$A\subset S$, its closed essential support is
\[
 \operatorname{ess\,supp}_S A
 =\{a\in S:\sigma(A\cap I)>0\text{ for every open arc }I\ni a\}.
\]
It is empty exactly when $A$ is null.  Almost every point of $A$ is
both a density-one direction for $A$ and a member of this support.
Changing the representative of a state on a null set changes neither
its support nor any assertion below.

\begin{lemma}[Strict state separation]
\label{lem:transitions-strict-separation}
Suppose $A,B\subset S$ are states at $x$ representing disjoint label
groups whose union is $\mathcal I\setminus\{X\}$, for a fixed label
$X$.  At any nonparallel density-one directions $a$ of $A$ and $b$ of
$B$, the point $x+a+b$ has spatial density one of $X$.

If $\sigma(A)>0$, the closed essential support of $B$ contains no
pair at chordal distance one.  In particular, it has a strict angular
separation from its image under $R$ and under $R^{-1}$, with the
empty-support case interpreted vacuously.

The separation is uniform along a sequence of centers and states
$(x_n,A_n,B_n)$ with the same represented groups, provided
$\ind_{A_n}\to\ind_A$ in $L^1(S)$ for a set $A$ of positive measure.
More precisely, for all sufficiently large $n$, angular differences
between two points of $\operatorname{ess\,supp}_S B_n$ stay a fixed
positive distance from both $\kappa$ and $-\kappa$ modulo $2\pi$.
\end{lemma}

\begin{proof}
The density conclusion of Lemma~\ref{lem:palette-transfer}, following
\eqref{eq:source-10}, gives density zero of each of the four represented
labels at the nonparallel sum.  Their complement, the singleton $X$,
therefore has density one there.

We prove the uniform assertion by contradiction; the fixed-center
case follows by using a constant sequence.  If it fails, choose pairs
in the supports of $B_n$ with chordal distances tending to one.  Each
support point is approximable by density-one directions of $B_n$.
We may thus choose density directions $b_n,d_n$ in those states such
that, after a subsequence,
\[
 b_n\to b,\qquad d_n\to d,\qquad
 w_n=b_n-d_n\to w=b-d,\qquad |w|=1.
\]
Choose a small angular interval $I$ with $\sigma(A\cap I)>0$, whose
closure is separated from the directions parallel to any of $w,b,d$.
Such an interval exists because only finitely many directions are
excluded and $A$ has positive measure.

For $a\in I$, the implicit function theorem gives a unit direction
$c_n(a)$, smoothly depending on $a$, with
\[
 |w_n+a-c_n(a)|=1,
 \qquad c_n\longrightarrow\id\quad\text{in }C^1(I).
\]
Indeed, at $w_n=w$ and $c_n(a)=a$ the squared-length equation is
satisfied and its derivative in the angle of $c_n(a)$ is
$-2w\cdot Ja\ne0$.  Shrinking $I$ around a density point of $A$ if
necessary makes this construction uniform on its closure.

The maps $c_n$ have uniformly bounded inverse derivatives.  It follows
that $\ind_{A_n}\circ c_n\to\ind_A$ in $L^1(I)$: the difference
between $\ind_{A_n}\circ c_n$ and $\ind_A\circ c_n$ is controlled by
the $L^1$ convergence under change of variables; continuous
approximation of $\ind_A$ and $c_n\to\id$ treat the remaining
difference.  Together with $\ind_{A_n}\to\ind_A$ on $I$, this shows
that the set of $a\in I$ for which both $a$ and $c_n(a)$ lie in $A_n$
has positive measure for large $n$.  Excluding null sets, choose such
an $a$ for which both are density-one directions.  The choice of $I$
ensures that $(a,b_n)$ and $(c_n(a),d_n)$ are nonparallel.  The points
\[
 x_n+a+b_n,\qquad x_n+c_n(a)+d_n
\]
both have density one of color $X$, yet their distance is
$|w_n+a-c_n(a)|=1$.  This contradicts properness of density-one points.
At a single center, compactness of the support turns exclusion of
distance-one pairs into the asserted strict margin.
\end{proof}

\begin{lemma}[Singleton--pair alternation is impossible]
\label{lem:transitions-two-state}
At a fixed center, there is no measurable state function with only the
two values $X,D$, representing respectively a singleton $\{X\}$ and
a fixed disjoint pair $D$, and changing state under $R$ almost
everywhere.
\end{lemma}

\begin{proof}
Put
\[
 \gamma=2\arcsin\frac1{2\sqrt3},
 \qquad \cos\gamma=\frac56.
\]
The ratio $\gamma/\pi$ is irrational.  Otherwise $e^{i\gamma}$ would
be a root of unity, making $2\cos\gamma=5/3$ a rational algebraic
integer, hence an integer, a contradiction.

There is a positive-measure set of angles $t$ whose states at $t$ and
$t+\gamma$ are equal.  If not, the state would flip almost everywhere
under $\gamma$, and hence be invariant under $2\gamma$.  A measurable
function invariant under an irrational circle rotation is constant
almost everywhere: its nonconstant Fourier coefficients vanish, since
the corresponding rotation eigenvalues differ from one.  This would
contradict the flip under $\kappa$.

Choose $t$ from the positive-measure equality set, also imposing all
finitely many needed $R$-alternation identities and density-direction
conditions.  For $j=0,1,2$, form the points
\[
 z_j=x+e^{i(t+2j\kappa)}+e^{i(t+\pi+\gamma+2j\kappa)}
     =x+e^{i(t+2j\kappa)}(1-e^{i\gamma}).
\]
They have distance $|1-e^{i\gamma}|=1/\sqrt3$ from $x$ and form an
equilateral triangle of side one.  In each sum the two nonparallel
endpoints have opposite states: $2\kappa$ preserves the state and
$\pi=3\kappa$ flips it.  The density consequence of
Lemma~\ref{lem:palette-transfer} therefore gives density one at each
$z_j$ to the union of the two labels outside $\{X\}\cup D$.
A uniform common translation in a sufficiently small disk puts all
three vertices in that union and in the typical set with positive
probability; the three failure probabilities tend to zero by density
one.  The resulting unit equilateral triangle would be properly
colored with two labels, a contradiction.
\end{proof}

\subsection{The cyclic-center locus}

\begin{theorem}[Local finiteness of cyclic centers]
\label{thm:cyclic-local-finiteness}
The set
\begin{equation}
 \mathcal L=\{x\in\R^2:\Gamma_x\text{ contains a cycle}\}
 \label{eq:source-19}
\end{equation}
is closed and locally finite: it has finite intersection with every
compact subset of the plane.
\end{theorem}

For any specified simple label cycle, its locus is the intersection
of its finitely many closed edge-presence loci, hence is closed.
There are only finitely many simple cycles on five labels.  It
therefore suffices to prove local finiteness separately for each
triangle, four-cycle, and five-cycle.

For a triangle, \eqref{eq:source-15} for its three edges and
Proposition~\ref{prop:palette-size} force $P_x(e)$ to intersect the
pair of outside labels almost everywhere.  Indeed, a nonempty subset
of the three triangle labels of size at most two lies in some
triangle edge, which is forbidden.  By \eqref{eq:source-12} the
intersection with the outside pair is unique almost everywhere and
alternates under $R$.  Thus every such $x$ is binary for that fixed
pair.  Lemma~\ref{lem:binary-separation} proves local finiteness of
the triangle locus.

\subsubsection{Four-cycles: continuity of states}

Fix a four-cycle and relabel its vertices $0,1,2,3$ in cyclic order,
with fifth label $X$.  Let $L$ be its closed locus and write
\[
 U=\{0,2\},\qquad V=\{1,3\}.
\]
For $x\in L$ define the canonical state $w_x(e)$ as follows: it is
$X$ if $X\in P_x(e)$, and otherwise it is $U$ or $V$ according to
which of these two pairs equals $P_x(e)$.  This is exhaustive almost
everywhere.  A palette not containing $X$ has size at most two, and
\eqref{eq:source-15} excludes every singleton and every adjacent pair
of the four-cycle, leaving precisely the two diagonals.  The states
represent the groups $\{X\},U,V$ contained in their palettes, and
\eqref{eq:source-12} gives
\[
 w_x(e)\ne w_x(Re)\quad\text{almost everywhere}.
\]
Both diagonal states have positive angular measure.  If either were
null, the other and $X$ would give the impossible alternation of
Lemma~\ref{lem:transitions-two-state}.

\begin{lemma}[Continuity of four-cycle states]
\label{lem:transitions-four-continuity}
For each $s\in\{X,U,V\}$, the map
$x\mapsto\ind_{\{w_x=s\}}$ is continuous from $L$ into $L^1(S)$.
\end{lemma}

\begin{proof}
At each $y\in L$ choose an avoiding trace $b_y^{ij}$ for each of its
four specified edges.  If $k$ belongs to the group represented by a
state $s$, then
\begin{equation}
 \ind_{\{w_y=s\}}
    \le\sum_{ij\text{ in the four-cycle}} b^{ij}_{y,k}
 \quad\text{almost everywhere}.
 \label{eq:source-20}
\end{equation}
To see this, fix a generic direction in state $s$.  If the palette
has another member besides $k$, some cycle edge contains that member
but does not contain $k$.  For a diagonal state, choose an edge
incident to its other member; for $k=X$, choose an edge incident to
the possible additional label.  If no other member exists, any
cycle edge suffices.  The corresponding trace avoids its edge
endpoints and is supported within the palette, so its $k$-coordinate
equals one.  There are only finitely many chosen traces, allowing one
common conull set of directions for this argument.

Let $y_n\to x$ in $L$.  Along any subsequence, pass further to weak-*
limits of the three state indicators and all four chosen traces.
Denote the former limits by $f_X,f_U,f_V$.  The trace limits belong
to $\mathcal W_x$ by its closed-graph property.  Inequality
\eqref{eq:source-20} passes to the weak-* limits, and its nonnegative
right-hand side shows that wherever $f_s>0$, every label represented
by $s$ belongs to $P_x$, almost everywhere.  Two different state
groups would require at least three palette labels, contradicting
Proposition~\ref{prop:palette-size}.  Since the $f_s$ are nonnegative
and sum to one, exactly one equals one at almost every direction.
Its group determines precisely the canonical state $w_x$: a diagonal
already fills both palette slots, while a positive singleton state
requires $X\in P_x$.  Thus $f_s=\ind_{\{w_x=s\}}$ almost everywhere.

Weak-* convergence of indicators to an indicator implies $L^1$
convergence: for indicators $f_n,f$, test against $1$ and $f$ in
$\int|f_n-f|=\int f_n+\int f-2\int f_nf$.  Every subsequence has a
further subsequence with this same $L^1$ limit, which proves the
claimed continuity.
\end{proof}

Let $S_{U,x},S_{V,x}$ be the closed essential supports of the two
diagonal states, and put $N_x=S_{U,x}\cup S_{V,x}$.  Both supports
are nonempty.  Across their nonparallel density directions, the
four represented labels force singleton density one of $X$.
Lemma~\ref{lem:transitions-strict-separation} therefore gives a strict
margin between each support and its own $R$-shift.  These margins
are uniform along any convergent sequence in $L$, by
Lemma~\ref{lem:transitions-four-continuity} and positivity of both
limiting diagonal states.  We also have the following consequence
of that continuity:
\[
 b\in S_{s,x},\quad y_n\to x\text{ in }L
 \quad\Longrightarrow\quad
 \text{there are }b_n\in S_{s,y_n}\text{ with }b_n\to b
 \quad(s=U,V).
\]
Indeed, every neighborhood of $b$ has positive limiting state
measure, hence positive state measure at all sufficiently large
$n$.  Successively smaller neighborhoods give the sequence.

The set $N_x$ is a proper subset of $S$.  Otherwise a point in the
intersection of the two supports would rotate under $R$ outside
both, by their strict separations.  The two nonempty supports would
therefore have to be disjoint and cover the connected circle, which
is impossible.  On every component gap of $S\setminus N_x$ the
state is $X$ almost everywhere.  Such a gap has angular length at
most $\kappa$, since an interval of greater length overlaps its
$R$-translate in an open interval, contrary to the state-change
rule.

\subsubsection{Four-cycles: transport and the accumulation contradiction}

For nonparallel $a,b\in S$ and $h$ sufficiently small, the local
inverse of the sum map defines
\[
 T_h(a,b)=(c,d),\qquad c+d=a+b-h,
 \qquad (c,d)\text{ near }(a,b).
\]
It is smooth in the two angular variables and in $h$, and for fixed
$h$ it is a local diffeomorphism of the angular pair.  Its branch at
$h=0$ is the identity.

We will use the following support transport rule.  If $y=x+h\in L$
and $a\in S_{U,x}$, $b\in S_{V,x}$ are nonparallel, then
\[
 T_h(a,b)\in N_y\times N_y
\]
whenever the indicated nearby branch is defined.  In arbitrarily
small product neighborhoods of $(a,b)$ there is a positive-measure
set of pairs of density directions for $U,V$ at $x$.  The spatial
sums of every such nonparallel pair have density one of $X$.
The local diffeomorphism $T_h$ permits removal of all null sets
needed to make both output directions density directions of their
states at $y$ and members of the corresponding essential supports.
Neither output state can be $X$: the density consequence of
Lemma~\ref{lem:palette-transfer} at $y$ would then give density zero
of $X$ at the same spatial sum.  Both outputs consequently belong
to $N_y$.  Approximating the initial support pair by these generic
pairs and using closedness proves the rule.  It also holds with
$U,V$ interchanged, and in the reverse direction under $T_{-h}$.
All null sets here concern the two fixed centers in a single
application; no joint measurability in the center is used.

Suppose now that distinct points $y_n=x+h_n$ of $L$ tend to $x$.
After a subsequence, write
\[
 \delta_n=|h_n|>0,\qquad H_n=h_n/\delta_n\longrightarrow H\in S.
\]
Take any endpoint $a$ of a gap of $N_x$, first with the gap just to
its right in real angular coordinates.  It belongs to at least one
diagonal support, say $S_{U,x}$; the other case is symmetric.
Assume $H\cdot a\ne0$.

The first output angle of $T_h(a,b)$ has directional velocity
\[
 V_H(a;b)=-\frac{H\cdot b}{a\times b}
   =-H\cdot Ja-(H\cdot a)\cot(\arg b-\arg a),
 \qquad a\times b=\sin(\arg b-\arg a).
\]
Indeed, differentiating $c+d=a+b-h$ at $h=0$ in direction $H$ gives
$Ja\,\dot\theta+Jb\,\dot\phi=-H$; take the dot product with $b$.
Since $H\cdot a\ne0$, this velocity takes different values at
projectively different nonparallel directions $b$.  The positive
measure of the $V$ state allows choices $b,b'\in S_{V,x}$, neither
parallel to $a$, with $V_H(a;b)>V_H(a;b')$.

Set $c_n=T_{h_n}^{1}(a,b)$.  By support transport, $c_n\in N_{y_n}$,
and smoothness gives
\[
 \arg c_n=\arg a+\delta_n V_H(a;b)+o(\delta_n).
\]
If $c_n\in S_{U,y_n}$ along an infinite subsequence, choose
$b'_n\in S_{V,y_n}$ with $b'_n\to b'$.  Transporting
$(c_n,b'_n)$ back by $T_{-h_n}$ gives a first endpoint in $N_x$
whose angle is
\[
 \arg a+\delta_n\bigl(V_H(a;b)-V_H(a;b')\bigr)+o(\delta_n).
\]
This lies strictly inside the gap for large $n$, a contradiction.
To justify the error term without any convergence rate for $b'_n$,
let $\tau(h,\theta,\beta)$ be a real lift of the first output angle
on a fixed branch near $(0,\arg a,\arg b')$.  The identity
$\tau(0,\theta,\beta)=\theta$ gives, uniformly on a smaller
neighborhood,
\[
 \tau(h,\theta,\beta)
    =\theta+D_h\tau(0,\theta,\beta)h+O(|h|^2).
\]
Thus the change in the auxiliary angle $\beta$ affects only the
coefficient multiplied by $h$.  Continuity of that coefficient,
$c_n\to a$, $b'_n\to b'$, and $H_n\to H$ give the claimed
$o(\delta_n)$ remainder.

It follows that $c_n\in S_{V,y_n}$ for all large $n$.  The strict
separation for $U$ leaves a neighborhood of $Ra$ free of $S_{U,x}$.
Immediately to the right of $a$ the state is $X$ almost everywhere,
so immediately to the right of $Ra$ it is non-$X$ almost everywhere.
It must therefore be $V$ on a one-sided neighborhood of $Ra$, giving
$Ra\in S_{V,x}$.  There are points of $S_{V,y_n}$ tending to $Ra$.
Together with $c_n\to a$, these violate the uniform strict separation
of $S_{V,y_n}$ from its own $R$-shift.

For a gap lying to the left of $a$, choose
$V_H(a;b)<V_H(a;b')$ and repeat the argument on that side.  Thus
every gap endpoint would have to satisfy $H\cdot a=0$.  There are
only two such directions, and they are antipodal.  A nonempty gap
has distinct endpoints and length at most $\kappa<\pi$, so its two
endpoints cannot both be these directions.  This contradiction
rules out accumulation in $L$ and proves local finiteness for the
specified four-cycle.

\subsubsection{Five-cycles: support lists and a global step ban}

Fix a five-cycle, labeled in cyclic order by $\Z/5\Z$, and a center
$x$ in its locus.  By \eqref{eq:source-15} and
Proposition~\ref{prop:palette-size}, almost every palette is exactly
one of its diagonals
\[
 D_i=\{i,i+2\},\qquad i\in\Z/5\Z.
\]
Write $w_x(e)=i$ for this state.  Two such diagonals are disjoint
exactly when their indices differ by one, so
\begin{equation}
 w_x(Re)=w_x(e)\pm1\pmod5
 \quad\text{almost everywhere}.
 \label{eq:source-21}
\end{equation}
Let $S_i$ be the closed essential support of state $i$, and let
\[
 S(a)=\{i:a\in S_i\}\qquad(a\in S)
\]
be its support list.  Every list is nonempty, since the finitely many
supports cover the circle: an uncovered direction would have a
neighborhood on which all five state sets were null.

Every nonempty $S_i$ has strict separation from its own $R$-shift.
Indeed, state $i$ then has positive measure, and
\eqref{eq:source-21} forces at least one of its adjacent states to
have positive measure.  Their two disjoint diagonals cover four
labels, so Lemma~\ref{lem:transitions-strict-separation} applies.
It follows that, for every direction $a$,
\[
 S(a)\cap S(Ra)=\varnothing.
\]
Also, each $i\in S(a)$ has an index adjacent to it in each of
$S(Ra)$ and $S(R^{-1}a)$.  Approach $a$ by generic directions of
state $i$, use \eqref{eq:source-21}, and pass to a constant adjacent
index along a subsequence and then to its closed support.

\begin{lemma}[A support contact bans an angular step]
\label{lem:transitions-step-ban}
If $i,i+1\in S(a)$ at any direction $a$, then the transitions
$i\to i+1$ and $i+1\to i$ under $R$ are absent almost everywhere
on the entire angular circle at $x$.
\end{lemma}

\begin{proof}
Set $z=x+a$.  There is a trace in $\mathcal W_z$ avoiding $D_i$.
In fact, in each neighborhood of $a$ the state-$i$ set has positive
measure.  By \eqref{eq:source-10}, the two-angle sum map, and Fubini,
almost every direction in this set has a typical actual unit-circle
sampling at the center $x$ plus that direction, avoiding $D_i$
almost everywhere.  Choose such directions tending to $a$ and pass
to a weak-* limit.  The same construction gives, separately, a trace
at $z$ avoiding $D_{i+1}$.  Consequently, almost everywhere,
\[
 P_z(f)\not\subset D_i,\qquad P_z(f)\not\subset D_{i+1}.
\]
On the other hand, \eqref{eq:source-11} applied to
$z+f=x+f+a$ gives
$P_z(f)\cap P_x(f)=\varnothing$ almost everywhere.  Its first
endpoint map is $f\mapsto f$, a local angular diffeomorphism; its
second endpoint is constant, which is permitted.  The excluded
parallel directions $f=\pm a$ form a null set.

Let $X$ be the label outside $D_i\cup D_{i+1}$.  Whenever $w_x(f)=i$,
disjointness from $P_x(f)=D_i$ leaves only $D_{i+1}\cup\{X\}$
available to $P_z(f)$, and the avoiding trace for $D_{i+1}$ forces
$X\in P_z(f)$.  The same reasoning applies when $w_x(f)=i+1$.
A step in either order between these states would therefore put
$X$ into both $P_z(f)$ and $P_z(Rf)$, contrary to
\eqref{eq:source-12}.  The statements used hold on a common conull
set, also after pulling back by $R$, proving the global step ban.
\end{proof}

\subsubsection{Five-cycles: the accumulation contradiction}

Suppose distinct centers $y_n=x+h_n$ of this same five-cycle locus
tend to $x$.  Pass to a subsequence with $h_n/|h_n|\to H\in S$.
For large $n$, write
\[
 y_n=x+a_n^++a_n^-,\qquad
 a_n^\pm=\frac{h_n}{2}
   \pm\sqrt{1-\frac{|h_n|^2}{4}}\,J\frac{h_n}{|h_n|}.
\]
These are unit directions and are nonparallel for each large $n$,
although their limits are the antipodal pair $a=JH$ and $-a$.

For either endpoint $s=a_n^+$ or $s=a_n^-$, the diagonals in $S(s)$
have no common label.  Suppose instead that every $D_i$ with
$i\in S(s)$ contained a label $k$.  The finitely many closed
supports not containing $s$ can be excluded by a sufficiently small
neighborhood of $s$.  Almost every direction in that neighborhood
therefore has $k\in P_x$.  In a product neighborhood of the
nonparallel pair $(a_n^+,a_n^-)$, \eqref{eq:source-10} and the local
diffeomorphism of the sum map now imply that $A_k$ has zero area
in an open spatial neighborhood of $y_n$.  This contradicts
$y_n\in H_k$, which holds for every label because all five cycle
edges are present at $y_n$.  The neighborhoods in this argument may
depend on $n$; no uniform nonparallelness of the endpoint pairs is
needed.

The no-common-label property passes to $a$ and $-a$.  In fact, for
any fixed direction $u$, finiteness and closedness of the supports
give $S(v)\subset S(u)$ for every $v$ in some neighborhood of $u$.
Thus eventually $S(a_n^+)\subset S(a)$, and intersecting the
diagonals over the larger list $S(a)$ still gives the empty set;
likewise at $-a$.  A nonempty index list whose diagonals have empty
intersection contains adjacent indices.  Otherwise it is an
independent set in the index five-cycle, of size at most two, and
any two nonadjacent diagonals have a common label.

Let $S^{(r)}=S(R^ra)$ for $r\in\Z/6\Z$.  Both $S^{(0)}$ and
$S^{(3)}$ contain adjacent index pairs.  Rotate the label indices
so that $0,1\in S^{(0)}$.  Then
\[
 2,4\in S^{(1)}\cap S^{(5)}:
\]
these lists must avoid $0,1$ and must contain a neighbor of each,
forcing $4$ for $0$ and $2$ for $1$.  If $j,j+1\in S^{(3)}$, the
same argument gives $j-1,j+2\in S^{(2)}\cap S^{(4)}$.  Adjacent
position lists are disjoint, so neither of $j-1,j+2$ can be $2$ or
$4$.  Checking the five possible residues of $j$ leaves only
$j=1$ or $j=4$: their pairs $\{j-1,j+2\}$ are respectively
$\{0,3\}$ and $\{3,1\}$; each of $j=0,2,3$ contains $2$ or $4$.

If $j=1$, the adjacent pairs at $a$ and $-a$ are $\{0,1\}$ and
$\{1,2\}$, so Lemma~\ref{lem:transitions-step-ban} bans both
possible steps from state $1$.  If $j=4$, the pairs are
$\{0,1\}$ and $\{4,0\}$, banning both steps from state $0$.
In either case \eqref{eq:source-21} makes the shared state null,
contradicting its membership in $S^{(0)}$.  Only the step bans at
the two fixed directions $a,-a$ have been used, so no union of
exceptional null sets over arbitrary directions is involved.
This rules out accumulation in the five-cycle locus.

Each specified simple-cycle locus is closed and has no accumulation
point.  Its intersection with a compact set is therefore finite,
by compactness.  Taking the finite union over all specified simple
cycles proves Theorem~\ref{thm:cyclic-local-finiteness}.
